%-------------------------------------------------------------------------------
% LATEX TEMPLATE PRESENTATIE
%-------------------------------------------------------------------------------
% Dit template is voor gebruik door studenten van de de bacheloropleiding 
% Informatica van de Universiteit van Amsterdam.
% Voor informatie over presenteren, zie 
%                               https://practicumav.nl/presenteren/index.html
%
%-------------------------------------------------------------------------------
%	PACKAGES EN CONFIGURATIE
%-------------------------------------------------------------------------------
% Gebruik de optie "sidebar" voor toevoeging van een sidebar met inhoudsopgave
% en de optie "dyslexic" voor gebruik van het OpenDyslexic-lettertype
\documentclass[aspectratio=169,sidebar,dyslexic]{uva-inf-presentation}
\usepackage[dutch]{babel}

% Vul hieronder de gevraagde gegevens in, 
% meerdere auteurs en UvAnetID's gescheiden door een puntkomma 
\title{Titel} 
\authors{Auteur 1; Auteur 2}
\uvanetids{UvAnetID student 1; UvAnetID student 2}
\course{}
\tutor{}
\docent{} 
\group{}
\programme{Informatica}

\begin{document}
%-------------------------------------------------------------------------------
%	AUTOMATISCHE SLIDES
%-------------------------------------------------------------------------------
\begin{titelframe}
\titlepage
\end{titelframe}

\begin{titelframe}
\frametitle{Inhoudsopgave}
\tableofcontents
\end{titelframe}

%-------------------------------------------------------------------------------
%	PRESENTATIE SLIDES
%-------------------------------------------------------------------------------
\section{Eerste sectie}
\subsection{Eerste subsectie}

\begin{frame}
\frametitle{Titel eerste slide}
\end{frame}

%-------------------------------------------------------------------------------
\end{document}
